%%----------------------------------------------------------------------------%%
%% Optional accessible design choices
%%----------------------------------------------------------------------------%%
%% This file contains optional visual choices only.
%%
%% Required accessibility settings, PDF metadata, hyperlink setup, MathML
%% support, ETD formatting, and student-specific macros belong elsewhere.
%%
%% This file is loaded early in dissertation-main-accessible.tex so that font
%% choices happen immediately after unicode-math. Visual choices are defined
%% here but applied later through \NCSUApplyOptionalVisualChoices.
%%----------------------------------------------------------------------------%%


%% -------------------------------------------------------------------------- %%
%% Optional: thesis font choices
%% -------------------------------------------------------------------------- %%
%% This template is designed for LuaLaTeX, so font changes should use fontspec
%% and unicode-math rather than older LaTeX font packages.
%%
%% Activate at most one font option near the bottom of this section.
%%
%% Leave all font options commented out to use the template's default font setup.

%% -------------------------------------------------------------------------- %%
%% Option 1: Libertinus
%% A polished, readable serif font with a matching math font.
%% -------------------------------------------------------------------------- %%

\newcommand{\NCSUUseLibertinusFonts}{%
  \setmainfont{Libertinus Serif}[
    Ligatures=TeX
  ]%
  \setsansfont{Libertinus Sans}[
    Ligatures=TeX,
    Scale=MatchLowercase
  ]%
  \setmonofont{Libertinus Mono}[
    Scale=MatchLowercase
  ]%
  \setmathfont{Libertinus Math}%
}

%% -------------------------------------------------------------------------- %%
%% Option 2: STIX Two
%% A formal, journal-like font family with strong math support.
%% -------------------------------------------------------------------------- %%

\newcommand{\NCSUUseSTIXTwoFonts}{%
  \setmainfont{STIX Two Text}[
    Ligatures=TeX
  ]%
  \setsansfont{TeX Gyre Heros}[
    Ligatures=TeX,
    Scale=MatchLowercase
  ]%
  \setmonofont{Latin Modern Mono}[
    Scale=MatchLowercase
  ]%
  \setmathfont{STIX Two Math}%
}

%% -------------------------------------------------------------------------- %%
%% Option 3: Latin Modern
%% This is close to the default LaTeX look and is the safest fallback.
%% -------------------------------------------------------------------------- %%

\newcommand{\NCSUUseLatinModernFonts}{%
  \setmainfont{Latin Modern Roman}[
    Ligatures=TeX
  ]%
  \setsansfont{Latin Modern Sans}[
    Ligatures=TeX,
    Scale=MatchLowercase
  ]%
  \setmonofont{Latin Modern Mono}[
    Scale=MatchLowercase
  ]%
  \setmathfont{Latin Modern Math}%
}

%% Activate at most one font option.
%% For the default template font, comment out all three lines.

\NCSUUseLibertinusFonts
% \NCSUUseSTIXTwoFonts
% \NCSUUseLatinModernFonts


%% -------------------------------------------------------------------------- %%
%% Optional: cleaner NC State numbered chapter headings
%% -------------------------------------------------------------------------- %%
%% This changes the visual appearance of numbered chapter headings without
%% changing the actual \chapter command. It avoids decorative chapter-heading
%% packages.
%%
%% It deliberately leaves \chapter* headings unchanged. Therefore, preliminary
%% headings such as ABSTRACT, DEDICATION, BIOGRAPHY, ACKNOWLEDGEMENTS, TABLE OF
%% CONTENTS, LIST OF TABLES, and LIST OF FIGURES retain the official formatting
%% defined in ncsuthesis.cls. Unnumbered headings such as REFERENCES and the
%% APPENDICES divider page also retain the official class formatting.
%%
%% Numbered appendix headings still receive this style because LaTeX implements
%% them as numbered chapters after \appendix.
%%
%% This option is disabled by default so that the official chapter-heading
%% formatting in ncsuthesis.cls remains authoritative. Turn it on only when
%% a deliberately different visual style is desired.

\newif\ifNCSUUseOptionalChapterHeadings

\NCSUUseOptionalChapterHeadingstrue
%\NCSUUseOptionalChapterHeadingsfalse

\makeatletter

\newcommand{\NCSUChapterTitleWithRule}[1]{%
  #1\par\nobreak
  \vspace{0.18in}%
  {\color{black}\rule{0.45\textwidth}{0.8pt}}%
}

\newcommand{\NCSUUseAccessibleChapterHeadings}{%
  \@ifundefined{ParseLaTeXeHeading}{%
    %% Traditional LaTeX chapter driver.
    \renewcommand{\@makechapterhead}[1]{%
      \vspace*{1.0in}%
      {\parindent \z@ \raggedright \normalfont
        \ifnum \c@secnumdepth >\m@ne
          {\normalsize\bfseries\color{black}%
          \MakeUppercase{\@chapapp}\space \thechapter\par}%
          \vspace{0.25in}%
        \fi
        {\Large\bfseries ##1\par}%
        \vspace{0.18in}%
        {\color{black}\rule{0.45\textwidth}{0.8pt}\par}%
        \vspace{0.75in}%
      }%
    }%
  }{%
    %% LaTeX 2026 heading-template instance used by tagged-PDF builds.
    \EditInstance{heading}{ncsu-chapter}{%
      after-penalty-vspace = 1.0in,
      after-vspace         = 0.75in,
      prefix               = \MakeUppercase{\@chapapp},
      heading-decls        = \raggedright\parindent\z@\normalfont,
      prefix-decls         = \normalsize\bfseries\color{black},
      number-decls         = \normalsize\bfseries\color{black},
      title-decls          = \Large\bfseries
    }%
    \EditInstance{headformat}{ncsu-chapter}{%
      heading-indent    = 0pt,
      prefix-number-sep = \WordSpaceAmount{1},
      number-title-sep  = 0.25in,
      title-format      = {\NCSUChapterTitleWithRule{##1}}%
    }%
  }%
}

\makeatother


%% -------------------------------------------------------------------------- %%
%% Optional: pretty theorem cards
%% -------------------------------------------------------------------------- %%
%% The required file defines theorem-like environments in a plain,
%% final-ETD-safe way.
%%
%% This optional block changes the appearance of those environments to a softer
%% math-textbook style. It does not use amsthm or tcolorbox.
%%
%% Turn this on or off with the toggle below.

\newif\ifNCSUPrettyTheorems

% \NCSUPrettyTheoremstrue
\NCSUPrettyTheoremsfalse


%% -------------------------------------------------------------------------- %%
%% Apply optional visual choices
%% -------------------------------------------------------------------------- %%
%% This command is called from dissertation-main-accessible.tex after
%% ncsu-required.tex has been loaded.
%%
%% Do not call this command inside optional-accessible.tex. It must run after
%% ncsu-required.tex because the chapter heading style uses NC State colors and
%% the theorem-card style uses the theorem macros defined there.

\makeatletter

\newcommand{\NCSUApplyOptionalVisualChoices}{%
  \ifNCSUUseOptionalChapterHeadings
    \NCSUUseAccessibleChapterHeadings
  \fi

  \ifNCSUPrettyTheorems
    \@ifundefined{NCSUStatementBegin}{%
      \PackageWarningNoLine{ncsu-template}{%
        Pretty theorem cards were requested, but the base theorem macros were not found.
        Make sure ncsu-required.tex is loaded before optional visual choices are applied.
      }%
    }{%
      \renewcommand{\NCSUStatementBegin}[4]{%
        \refstepcounter{ncsuStatement}%
        \def\NCSUCurrentStatementRuleColor{##4}%
        \par\addvspace{12pt}%
        \noindent
        \begingroup
          \setlength{\fboxsep}{6pt}%
          \colorbox{##3}{%
            \makebox[\dimexpr\linewidth-2\fboxsep\relax][l]{%
              \strut\NCSUStatementTitle{##1}{\thencsuStatement}{##2}%
            }%
          }%
        \endgroup
        \par\nobreak\vspace{6pt}%
        \begingroup
          \leftskip=1em
          \rightskip=1em
          \noindent\ignorespaces
      }%

      \renewcommand{\NCSUStatementEnd}{%
        \par
        \endgroup
        \vspace{4pt}%
        \noindent{\color{\NCSUCurrentStatementRuleColor}\rule{\linewidth}{0.7pt}}%
        \par\addvspace{12pt}%
      }%
    }%
  \fi
}

\makeatother